\section{Introduction}
\label{sec:intro}

The deployment of massive, multi-task deep neural networks on resource-constrained edge devices has catalyzed a shift from dense monolithic architectures toward modular, parameter-efficient serving paradigms. Prominent among these are Mixture-of-Experts (MoE) models \cite{shazeer2017outrageously, moe_gshard, switch_transformer} and dynamic adapter-merging systems \cite{lora_hu, task_arithmetic_ilharco, dynamic_adapter_selection}, which dynamically route activations or average pre-trained weight matrices to adapt to incoming queries. On-device edge serving, however, is rarely homogeneous. Instead, a serving system must typically process a continuous, sequential stream of heterogeneous user requests, where the underlying task signature can switch abruptly from sample to sample (e.g., from creative writing to formal coding, then to translation).

In this sequential serving setting, existing routing controllers suffer from a fundamental, unresolved \emph{accuracy-stability dilemma}. This dilemma manifests as two pathological extremes:
\begin{itemize}
    \item \textbf{High-Frequency Spatial Jitter (Stateless Routers):} Traditional stateless routing controllers (e.g., SABLE \cite{sable_2025} and SPS-ZCA \cite{sps_zca_2025}) compute ensembling weights at each layer independently based on local representations or centroids. Under noisy or slightly out-of-distribution queries, this locality causes the routing coefficients to oscillate violently across network depth (layer-to-layer), a phenomenon known as the \emph{routing jitter paradox}. This spatial oscillation triggers severe representation drift, leading to downstream cascading representation collapse and degraded output quality.
    \item \textbf{Temporal Serv-Time Lag and Hysteresis (Stateful Routers):} To filter out spatial jitter, stateful systems (such as ChemMerge \cite{chemmerge_2026} or PAC-Kinetics \cite{pac_kinetics_2026}) introduce continuous-time biochemical kinetics or differential equations that act as a low-pass temporal filter. However, by maintaining an internal historical state across sequential samples, these models introduce an artificial "inertia." When a sudden task switch occurs in a heterogeneous query stream, these stateful controllers take several inference steps to "warm up" and adapt, exhibiting severe representational hysteresis. This lag results in catastrophic accuracy drops immediately following task transitions.
\end{itemize}

This trade-off raises a fundamental question: \emph{Is it possible to achieve globally smooth, jitter-free routing trajectories across network depth while maintaining zero temporal lag and zero serving hysteresis across samples?}

\begin{figure*}[t]
  \centering
  \begin{subfigure}[b]{0.33\textwidth}
    \centering
    \includegraphics[width=\textwidth]{results/fig1_routing_weights.png}
    \caption{Layer-wise Trajectory}
    \label{fig:layer_trajectory}
  \end{subfigure}
  \hfill
  \begin{subfigure}[b]{0.33\textwidth}
    \centering
    \includegraphics[width=\textwidth]{results/fig2_representational_lag.png}
    \caption{Hysteresis \& Temporal Lag}
    \label{fig:temp_lag}
  \end{subfigure}
  \hfill
  \begin{subfigure}[b]{0.33\textwidth}
    \centering
    \includegraphics[width=\textwidth]{results/fig3_pareto_frontier.png}
    \caption{Accuracy-Jitter Pareto}
    \label{fig:pareto_front}
  \end{subfigure}
  \caption{(a) Visualizing layer-wise routing weights under Homogeneous streams: SABLE oscillates violently, while ChemMerge and QPathMerge are smooth. (b) Weight adaptation after a task switch: Stateful kinetics (ChemMerge and PAC-Kinetics) exhibit severe lag, while SABLE and QPathMerge adapt instantly. (c) QPathMerge sweep of transition leakage $M$ maps a clear accuracy-jitter Pareto frontier.}
  \label{fig:main_results}
\end{figure*}

In this paper, we answer this question in the affirmative. Drawing a radical, visionary inspiration from statistical mechanics and path-integral formulations, we propose **Markovian Path-Integral Ensembling (QPathMerge)**. Rather than modeling routing as a feedforward heuristic or a continuous-time temporal filter, we view network depth as a discrete 1D lattice (representing sequential states) and model the entire routing trajectory as a discrete Euclidean path integral over the network layers. 

Specifically, we define the ensembling weight of each expert at each layer as the marginal probability of a path, where paths are governed by a Boltzmann distribution under a Euclidean action. The action contains potential energy (matching how well an expert fits a representation) and kinetic energy (penalizing expert switches between adjacent layers). This maps the layer-wise routing problem directly onto a 1D chain-structured Markov Random Field (MRF) \cite{markov_random_fields}. Rather than integrating slow differential equations, we execute the mathematically exact \textbf{Forward-Backward sum-product algorithm (Belief Propagation)} \cite{belief_prop_pearl} to resolve the global partition function and compute exact marginal ensembling weights for each layer in $O(L K^2)$ time.

Because QPathMerge propagates messages symmetrically forward and backward across the lattice depth for \emph{each input sample independently}, it achieves two profound advantages:
\begin{enumerate}
    \item \textbf{Decoupled Spatial Smoothing:} The transition barrier in the Euclidean action acts as a rigorous spatial low-pass filter, forcing adjacent layers to coordinate their expert selections and eliminating spatial jitter.
    \item \textbf{Absolute Statelessness (Zero Lag):} Since message propagation occurs entirely within the depth lattice of a single forward pass, the system carries zero state from one sample to the next. The temporal transition barrier is exactly zero. Under heterogeneous query streams, QPathMerge adapts instantly to task switches with zero lag and zero hysteresis.
\end{enumerate}

To make this path-integral formulation viable for low-power edge hardware, we also introduce a highly-efficient single-pass variant, \textbf{Recursive On-The-Fly QPathMerge (QPathMerge-Single)}, which we position as the primary deployment candidate for production serving. By executing a single forward pass and recursively updating backward messages over a Truncated Backward Horizon ($H=4$), QPathMerge-Single eliminates the dual-pass computational overhead entirely while preserving near-oracle spatial smoothness.

We evaluate both our theoretical bidirectional model and our single-pass deployment candidate in a high-fidelity 14-layer Coordinate Sandbox against seven state-of-the-art routing baselines, alongside a physical validation on a pre-trained ResNet-18 model using natural ImageNet-1K streams. Under rapid task-switching heterogeneous query streams, our framework slashes spatial layer-wise jitter by over $3\times$ compared to SABLE and ChemMerge while maintaining leading serving accuracy and zero temporal hysteresis, resolving the accuracy-stability serving trade-off.
